% GENERATED FILE. Edit canonical lessons, then run npm run generate:books.
% canonical-source-hash: fnv1a64:6c8ae0e2611a7b05
% canonical-lessons: FA-C507-akkas, FA-C507-akkasi, FA-C507-baghbani, FA-C507-ashpazi-kardan, FA-C507-koleksiyon

\chapter{Hobbies}
\label{ch:fa-hobbies}
\hlchaptermodality{507}

\begin{chapteropening}
\textbf{By the end of this chapter:} \emph{I can talk about my hobbies.}

Complete the last lesson of chapter 507: I can talk about my hobbies.
\end{chapteropening}

\section[\texorpdfstring{\fa{عکاس} ('akkâs)}{عکاس ('akkâs)}]{\fa{عکاس} ('akkâs) — a photographer}

\label{lesson:FA-C507-akkas}



\begin{quote}
Before the new one: say the Persian for to visit each other, then the Persian for reputation.
\end{quote}

\subsection*{You'll want to know: \fa{عکاس}}
\textbf{\fa{عکاس}} (\emph{'akkâs}) — \textquotedblleft{}a photographer\textquotedblright{}.

\begin{quote}
\fa{یک} \fa{عکاس} \fa{از} \fa{مسابقه} \fa{عکس} \fa{گرفت}.
\end{quote}

\begin{quote}
\emph{yek 'akkâs az mosâbeqe 'aks gereft.}
\end{quote}

\begin{quote}
A photographer took photos of the match.
\end{quote}

Built on \emph{'aks} 'photo'. \emph{'aks gereftan} — to take a photo.

\begin{cousinweb}
\emph{'akkâs} is built on \emph{'aks}, 'picture', a loan usually traced to Arabic \emph{ʿaks}, 'reflection', in the Arabic pattern for a person who does something.
\end{cousinweb}

\begin{grammarlens}[title={What you've built}]
One more word for telling, explaining and describing at B1.
\end{grammarlens}

\subsection*{Your turn}
\begin{itemize}
\raggedright
  \item \emph{Say it:} \emph{'akkâs}
  \item \emph{Say it:} \emph{'akkâs}, once more
  \item \emph{Say it:} a sentence of your own with \emph{'akkâs}
  \item \emph{Recall:} say the Persian for to visit each other, then the Persian for reputation, then say \emph{'akkâs} again
\end{itemize}

\begin{tcolorbox}[breakable,colback=teal!4,colframe=teal!35!black,title={Before you move on}]
What does \fa{عکاس} mean? (\textquotedblleft{}A photographer\textquotedblright{}.) Say it once more.
\end{tcolorbox}
\section[\texorpdfstring{\fa{عکاسی} ('akkâsi)}{عکاسی ('akkâsi)}]{\fa{عکاسی} ('akkâsi) — photography}

\label{lesson:FA-C507-akkasi}



\begin{quote}
Before the new one: say the Persian for reputation, then the Persian for a photographer.
\end{quote}

\subsection*{You'll want to know: \fa{عکاسی}}
\textbf{\fa{عکاسی}} (\emph{'akkâsi}) — \textquotedblleft{}photography\textquotedblright{}.

\begin{quote}
\fa{سرگرمی} \fa{پدرم} \fa{عکاسی} \fa{است}.
\end{quote}

\begin{quote}
\emph{sargarmi-ye pedaram 'akkâsi ast.}
\end{quote}

\begin{quote}
My father's hobby is photography.
\end{quote}

\emph{'akkâs} 'photographer' + \emph{-i}. \emph{'akkâsi kardan} — to take photographs.

\begin{cousinweb}
\emph{'akkâsi} is built on \emph{'akkâs}, 'photographer', with the Persian noun ending \emph{-i} (EtymDB); \emph{'akkâs} itself is built on \emph{'aks}, a loan from Arabic.
\end{cousinweb}

\begin{grammarlens}[title={What you've built}]
One more word for telling, explaining and describing at B1.
\end{grammarlens}

\subsection*{Your turn}
\begin{itemize}
\raggedright
  \item \emph{Say it:} \emph{'akkâsi}
  \item \emph{Say it:} \emph{'akkâsi}, once more
  \item \emph{Say it:} a sentence of your own with \emph{'akkâsi}
  \item \emph{Recall:} say the Persian for reputation, then the Persian for a photographer, then say \emph{'akkâsi} again
\end{itemize}

\begin{tcolorbox}[breakable,colback=teal!4,colframe=teal!35!black,title={Before you move on}]
What does \fa{عکاسی} mean? (\textquotedblleft{}Photography\textquotedblright{}.) Say it once more.
\end{tcolorbox}
\section[\texorpdfstring{\fa{باغبانی} (bâghbâni)}{باغبانی (bâghbâni)}]{\fa{باغبانی} (bâghbâni) — gardening}

\label{lesson:FA-C507-baghbani}



\begin{quote}
Before the new one: say the Persian for a photographer, then the Persian for photography.
\end{quote}

\subsection*{You'll want to know: \fa{باغبانی}}
\textbf{\fa{باغبانی}} (\emph{bâghbâni}) — \textquotedblleft{}gardening\textquotedblright{}.

\begin{quote}
\fa{مادربزرگ} \fa{در} \fa{بهار} \fa{باغبانی} \fa{میکند}.
\end{quote}

\begin{quote}
\emph{mâdarbozorg dar bahâr bâghbâni mikonad.}
\end{quote}

\begin{quote}
Grandmother gardens in spring.
\end{quote}

\emph{bâgh} 'garden' + \emph{-bân} 'keeper' + \emph{-i}. A gardener is \emph{bâghbân}.

\begin{cousinweb}
\emph{bâghbâni} is built on native Persian parts: \emph{bâgh}, 'garden', usually traced to Middle Persian \emph{bāg}, the ending \emph{-bân}, 'keeper', and \emph{-i}.
\end{cousinweb}

\begin{grammarlens}[title={What you've built}]
One more word for telling, explaining and describing at B1.
\end{grammarlens}

\subsection*{Your turn}
\begin{itemize}
\raggedright
  \item \emph{Say it:} \emph{bâghbâni}
  \item \emph{Say it:} \emph{bâghbâni}, once more
  \item \emph{Say it:} a sentence of your own with \emph{bâghbâni}
  \item \emph{Recall:} say the Persian for a photographer, then the Persian for photography, then say \emph{bâghbâni} again
\end{itemize}

\begin{tcolorbox}[breakable,colback=teal!4,colframe=teal!35!black,title={Before you move on}]
What does \fa{باغبانی} mean? (\textquotedblleft{}Gardening\textquotedblright{}.) Say it once more.
\end{tcolorbox}
\section[\texorpdfstring{\fa{آشپزی} \fa{کردن} (âshpazi kardan)}{آشپزی کردن (âshpazi kardan)}]{\fa{آشپزی} \fa{کردن} (âshpazi kardan) — to cook, to do the cooking}

\label{lesson:FA-C507-ashpazi-kardan}



\begin{quote}
Before the new one: say the Persian for photography, then the Persian for gardening.
\end{quote}

\subsection*{You'll want to know: \fa{آشپزی} \fa{کردن}}
\textbf{\fa{آشپزی} \fa{کردن}} (\emph{âshpazi kardan}) — \textquotedblleft{}to cook, to do the cooking\textquotedblright{}.

\begin{quote}
\fa{پدرم} \fa{آخر} \fa{هفته} \fa{ها} \fa{آشپزی} \fa{میکند}.
\end{quote}

\begin{quote}
\emph{pedaram âkhar-e haftehâ âshpazi mikonad.}
\end{quote}

\begin{quote}
My father cooks at weekends.
\end{quote}

\emph{âshpaz} 'cook' + \emph{-i} + \emph{kardan}. \emph{pokhtan} is to cook a particular dish.

\begin{cousinweb}
\emph{âshpazi} is built on native Persian parts: \emph{âsh}, 'soup', \emph{paz}, the present stem of \emph{pokhtan}, 'to cook', and \emph{-i}: soup-cooking. \emph{kardan} is inherited too.
\end{cousinweb}

\begin{grammarlens}[title={What you've built}]
One more word for telling, explaining and describing at B1.
\end{grammarlens}

\subsection*{Your turn}
\begin{itemize}
\raggedright
  \item \emph{Say it:} \emph{âshpazi kardan}
  \item \emph{Say it:} \emph{âshpazi kardan}, once more
  \item \emph{Say it:} a sentence of your own with \emph{âshpazi kardan}
  \item \emph{Recall:} say the Persian for photography, then the Persian for gardening, then say \emph{âshpazi kardan} again
\end{itemize}

\begin{tcolorbox}[breakable,colback=teal!4,colframe=teal!35!black,title={Before you move on}]
What does \fa{آشپزی} \fa{کردن} mean? (\textquotedblleft{}To cook, to do the cooking\textquotedblright{}.) Say it once more.
\end{tcolorbox}
\section[\texorpdfstring{\fa{کلکسیون} (koleksiyon)}{کلکسیون (koleksiyon)}]{\fa{کلکسیون} (koleksiyon) — a collection}

\label{lesson:FA-C507-koleksiyon}



\begin{quote}
Before the new one: say the Persian for gardening, then the Persian for to cook.
\end{quote}

\subsection*{You'll want to know: \fa{کلکسیون}}
\textbf{\fa{کلکسیون}} (\emph{koleksiyon}) — \textquotedblleft{}a collection\textquotedblright{}.

\begin{quote}
\fa{او} \fa{یک} \fa{کلکسیون} \fa{تمبر} \fa{قدیمی} \fa{دارد}.
\end{quote}

\begin{quote}
\emph{u yek koleksiyon-e tambr-e qadimi dârad.}
\end{quote}

\begin{quote}
He has a collection of old stamps.
\end{quote}

Plural \emph{koleksiyonhâ}.

\begin{cousinweb}
\emph{koleksiyon} is a modern loan, usually traced to French \emph{collection}. EtymDB has no entry for the Persian word.
\end{cousinweb}

\begin{grammarlens}[title={What you've built}]
One more word for telling, explaining and describing at B1.
\end{grammarlens}

\subsection*{Your turn}
\begin{itemize}
\raggedright
  \item \emph{Say it:} \emph{koleksiyon}
  \item \emph{Say it:} \emph{koleksiyon}, once more
  \item \emph{Say it:} a sentence of your own with \emph{koleksiyon}
  \item \emph{Recall:} say the Persian for gardening, then the Persian for to cook, then say \emph{koleksiyon} again
\end{itemize}

\begin{tcolorbox}[breakable,colback=teal!4,colframe=teal!35!black,title={Before you move on}]
What does \fa{کلکسیون} mean? (\textquotedblleft{}A collection\textquotedblright{}.) Say it once more.
\end{tcolorbox}
